% from ega0-22-fr.tex
\oldpage[0\textsubscript{IV-1}]{197}

vérifiant les conditions de (21.3.1). Soit en outre $\mathfrak m$ un idéal
\emph{de carré nul} dans $A$~; posons
\[
  K=A/\mathfrak m,\qquad
  B_{(K)}=B\otimes_AK=B/Bi(\mathfrak m)
\]
et soient
\[
  \varphi:A\longrightarrow K,\qquad
  \psi:B\longrightarrow B_{(K)}
\]
les homomorphismes canoniques.

Comme $\varphi$ est surjectif, $\Omega^1_{B_{(K)}/A}$ s'identifie
canoniquement à $\Omega^1_{B_{(K)}/K}$ (20.5.15). Notons que, puisque
$p\geqslant2$, pour tout $x\in\mathfrak m$, on a $j(i(x))=x^p=0$, autrement
dit $j(i(\mathfrak m))=0$~; par passage aux quotients on déduit donc de $j$ un
homomorphisme
\[
  j':B_{(K)}\longrightarrow A
\]
tel que, si $i'=\psi\circ i:A\longrightarrow B_{(K)}$, on ait
\[
  j'(i'(x))=x^p\qquad\text{et}\qquad i'(j'(y))=y^p
  \qquad\text{pour }x\in A\text{ et }y\in B_{(K)}.
\]
En posant, suivant les notations de (21.3.2),
\[
  \Theta_{B_{(K)}/A}
  =\Omega^1_{B_{(K)}/A}\otimes_{B_{(K)}}A_{[j']}
  =\Omega^1_{B_{(K)}/K}\otimes_{B_{(K)}}A_{[j']},
\]
on a donc (21.3.2.2) un homomorphisme canonique
\begin{equation}
\label{0.22.4.5.1-fr}
  \pi_{B_{(K)}/A}:\Theta_{B_{(K)}/A}\longrightarrow\Omega^1_A.
  \tag{22.4.5.1}
\end{equation}

D'autre part, on déduit de $i$ et $j$, par passage aux quotients (tenant compte
de ce que $j(i(\mathfrak m))=0$), des homomorphismes
\[
  K\xrightarrow{i_0}B_{(K)}\xrightarrow{j_0}K
\]
de sorte que l'on a le diagramme commutatif
\[
  \xymatrix{
    A \ar[r]^i \ar[d]_\varphi & B \ar[r]^j \ar[d]^\psi & A \ar[d]^\varphi \\
    K \ar[r]_{i_0} & B_{(K)} \ar[r]_{j_0} \ar[ur]_{j'} & K
  }
\]
Il est clair en outre que l'on a
\[
  j_0(i_0(s))=s^p\qquad\text{et}\qquad i_0(j_0(t))=t^p
  \qquad\text{pour }s\in K\text{ et }t\in B_{(K)}.
\]
On définit donc de même
\[
  \Theta_{B_{(K)}/K}
  =\Omega^1_{B_{(K)}/K}\otimes_{B_{(K)}}K_{[j_0]}
  =\Theta_{B_{(K)}/A}\otimes_AK,
\]
donc, en tensorisant (22.4.5.1) par $K$, on obtient un
$K$-\emph{homomorphisme canonique}
\begin{equation}
\label{0.22.4.5.2-fr}
  \pi'_{B_{(K)}/K}=\pi_{B_{(K)}/A}\otimes1_K:
  \Theta_{B_{(K)}/K}\longrightarrow\Omega^1_A\otimes_AK.
  \tag{22.4.5.2}
\end{equation}

